\documentclass[12pt]{article}

\usepackage{amssymb,amsmath,amsthm} % Include AMS Math symbols
\usepackage{graphicx}               % Allows for importation of figures 
\usepackage{float} % to get the table where i want it
\usepackage{booktabs} % professional table typesetting commands 

\usepackage{color}   % Sets up a nice typesetting environment for Python
\usepackage{listings}
\definecolor{dkgreen}{rgb}{0,0.6,0}
\definecolor{gray}{rgb}{0.5,0.5,0.5}
\definecolor{mauve}{rgb}{0.58,0,0.82}
\lstset{frame=tb,
  language=Python,
  aboveskip=3mm,
  belowskip=3mm,
  showstringspaces=false,
  columns=flexible,
  basicstyle={\small\ttfamily},
  numbers=left,
  numberstyle=\tiny\color{gray},
  keywordstyle=\color{blue},
  commentstyle=\color{dkgreen},
  stringstyle=\color{mauve},
  breaklines=true,
  breakatwhitespace=true,
  tabsize=3
}

% define command \expandableinput  for table input
\makeatletter
\def\expinput#1{\@@input#1 }
\makeatother

\title{ HW \# 1\\ Math/CS 471, Fall 2026}
\author{Theo Nichols \\ Andrew Vargas}
\date{\today}   % Activate to display a given date or no date


\begin{document}

% Display title
\maketitle
\clearpage

\section{Question 1: One Dimensional Newton's Method}

\subsection{Description}
We are tasked with finding the roots of three different scalar functions and determining the convergence rate, or rate at which our computational methods approach the root. Two of these are non-linear, and one is a simple linear function, the importance of which will be explored in subsequent sections. We are provided the Python code for and instructed to use Newton's Method, an iterative method which uses an initial guess for a root, or zero, of the function, and then repeats the formula below, where $x_{n+1}$ is the new guess for the root, $x_n$ is the previous, $f$ refers to the function, and $f'$ its derivative.
\begin{equation*}
    x_{n+1} = x_n - \frac{f(x_n)}{f'(x_n)}
\end{equation*}

\subsection{Method}
We used the provided code, and made some simple modifications to determine the convergence rate of Newton's Method, as well as place safeguards that ensure the code would iterate an appropriate number of times or halt if it would otherwise iterate indefinitely. To accomplish that, we simply added a user defined tolerance for the error estimate, and a maximum number of iterations that if reached would likely imply Newton's method was going to fail to find a root. Of course, this involved calculating an approximation of the error at each iteration. Our error estimate was calculated as $|x_{n+1} - x_n|$. \\ \\If this value  was found to be less than the user defined tolerance, the code would break. We were advised to choose $10^{-10}$ as our tolerance, and a maximum number of iterations that would give the method time to converge. I chose 100 rather arbitrarily, and this number of iterations was never reached. \\ \\We then calculated the convergence rate. The convergence rate refers to the rate at which a method approaches a true root. We tested for linear and quadratic convergence, where linear convergence implies that the error of the next iteration is proportional to a constant times the error of the current error, and quadratic convergence is more rapid as the error at the next step is proportional to the square of the current error.
To calculate the convergence rate, we took the ratio of the error estimate for the next iteration divided by that of the previous iteration. We first populated an array with these error estimates, then simply calculated $\frac{errors[i+1]}{errors[i]}$ and $\frac{errors[i+1]}{errors[i]^2}$ for linear and quadratic convergence respectively.\\ 
\\If a ratio remained nearly constant for every iteration, this implied the method was converging at that rate, as the errors were proportional. If on the other hand the values continued to change, this meant the errors were not proportional and the method was not converging at that rate. Below is a snippet of our Python code for computing these ratios. 
\begin{verbatim}
k_l = 0
k_q = 0
errors = []

for i in range(len(data) - 1):
    errors.append(abs(data[i+1] - data[i]))

for i in range(len(errors) -1):
    k_l = errors[i+1] / errors[i]
    k_q = errors[i+1] / (errors[i]**2)
        
    print("Linear convergence ration: " + str(k_l))
    print("Quadratic convergence ration: " + str(k_q))

\end{verbatim}
Data was an array containing all of the potential roots generated by the newton.py file provided to us. We piped the output of this driver code into a text file and the results can be seen in Table 1 below.

\begin{table}[H]
\centering
\begin{tabular}{c|c|c}
\toprule
Function & Linear Ratio & Quadratic Ratio \\
\midrule

$f(x)=x^2$ & $0.5$ & $4$ \\
         & $0.5$ & $8$ \\
         & $0.5$ & $16$ \\
         & $\vdots$ & $\vdots$ \\
         & $0.5$ & $2.684\times10^8$ \\
         & $0.5$ & $5.369\times10^8$ \\
         & $0.5$ & $1.074\times10^9$ \\
         & $0.5$ & $2.147\times10^9$ \\
         & $0.5$ & $4.295\times10^9$ \\

\midrule

$f(x)=x$ & -- & -- \\

\midrule

$f(x)=\sin(x)+\cos(x^2)$
& $0.06526547$ & $0.81112334$ \\
& $0.00405088$ & $0.77138257$ \\
& $1.633\times10^{-5}$ & $0.76775071$ \\
& $3.195\times10^{-7}$ & $919.708$ \\

\bottomrule
\end{tabular}
\caption{Functions and Their Computed Convergence Ratios}
\label{tab:1}
\end{table}

\subsection{Interpretation} 
As can be seen in Table 1, the linear ratio for $f(x) = x^2$ remained a constant 0.5. This indicates that Newton's did indeed converge linearly for this particular function and initial guess. The reason this occurs is due to the behavior of $f'(x)$ near the root. \\\\If we rewrite Newton's Method, $x_{n+1} = x_n - \frac{f(x_n)}{f'(x_n)}$ using the function \\$f(x) = x^2$, we get $x_n - \frac{x^2_n}{2x_n} = x_n - \frac{1}{2}x_n = \frac{1}{2}x_n$ which is precisely the definition of linear convergence, as the error is decreasing by a constant factor at each step. The root encountered here is a multiple root, that is, the function touches zero here but then rebounds in the positive y direction. The quadratic ratio continues to grow, indicating that the error is not proportional to the square of the error at each step. \\\\For $f(x) = x$, the reason there is no data is due to the fact that it is a linear function. As a result, Newton's method converges in one step, that is, with an initial guess of -0.5, the first computation yields \\$f(-0.5) - \frac{f(-0.5)}{f'(-0.5)} = -0.5 - -0.5 = 0$. As such there is no error, and the actual root of $x = 0$ has been found in a single step. \\ \\For $f(x) = sin(x) + cos(x^2)$, the linear ratio does not remain constant, but the quadratic ratio remains nearly constant, thus indicating quadratic convergence. The anomalous last value occurs because as Newton's Method gets closer to the root, the error becomes very small. Since the quadratic ratio divides the next error by the square of the current error, squaring an already very small error produces an extremely small denominator, causing the ratio to become very large. This does not indicate that the method is not converging quadratically, but is the result of a lack of numerical precision on my machine.

\section{Question 2: Multi-Dimensional Newton's Method}

Repeat same report structure for all following questions\\\\


\section*{Useful Latex Examples To Get You Started}

\subsection{Equations}
It has very powerful support for equations such as these:
\begin{equation}
f(x) = e^{1/2 - \sin(5 \pi    x)}, \label{eq:overcomplicated_exp_function} 
\end{equation}
\begin{equation}
p(x) = 1+x+x^2. \label{eq:special_polynomial} 
\end{equation}
% The label command is used to provide a unique identifier for each equation that is independent of any numbering

\noindent
Here is how to typeset fractions and square-roots
\begin{equation}
    r = \frac{-b + \sqrt{b^2 - 4ac}}{2a}
\end{equation}

\noindent
For inline equation, you would use the dollar sign, i.e., $p(x) = 1 + x + x^2$.

\noindent
If you want an un-numbered equation use the asterisk
\begin{equation*}
    q(x) = x^3
\end{equation*}

\noindent
If you want a multi-line equation use ampersands for the alignment column and then two backslashes to denote a line break.
\begin{align}
   p(x) & = (x + 1)(x + 2)(x) \\
        & = (x^2 + 3x + 2)(x) \\
        & = x^3 + 3x^2 + 2x
\end{align}

\subsection{Referencing}
\LaTeX also supports a very convenient labeling and referencing system
through the command \texttt{\\label} and \texttt{\\ref}.  Above, equation
(\ref{eq:special_polynomial}) is a polynomial and the function described by equation
(\ref{eq:overcomplicated_exp_function}) is plotted in Figure \ref{fig:1}. 

{\bf Note} that you may have to compile the document two times to get
all the cross-referencing between labels to work correctly.

\subsection{Figures}
Figures are also easy to incorporate. It is often best to use a figure
format that is based on vector graphics, e.g. Encapsulated PostScript
(eps). PNG is another good format.  An example is given in Figure \ref{fig:1}. 

% \begin{figure}[htb]
% \begin{center}
% \includegraphics[width=0.75\textwidth]{sample.png}
%    \caption{This figure displays the function from equation (\ref{eq:overcomplicated_exp_function}).}
%     \label{fig:1}
% \end{center}
% \end{figure}

\subsection{Tables}
Tables are often useful to display data more compactly than figures.


\subsection{Code Typesetting}

Simple typesetting
\begin{verbatim}
k_l = 0
k_q = 0
errors = []

for i in range(len(data) - 1):
    errors.append(abs(data[i+1] - data[i]))

for i in range(len(errors) -1):
    k_l = errors[i+1] / errors[i]
    k_q = errors[i+1] / (errors[i]**2)
        
    print("Linear convergence ration: " + str(k_l))
    print("Quadratic convergence ration: " + str(k_q))

\end{verbatim}


Fancy typesetting for code
\begin{lstlisting}
from scipy import *
from matplotlib import pyplot as plt
a = 4
b = 2
c = a*b
if c > 2:
    x = linspace(0,1,100)
    y = sin(x)
    plt.plot(x,y)
    plt.show()
\end{lstlisting}

%\bibliographystyle{plain}
%\bibliography{bibfile}

\end{document}  
